\documentclass[sigconf,anonymous,review]{acmart}
\settopmatter{printacmref=true}
\usepackage{booktabs}
\usepackage{amsmath}
\usepackage{graphicx}
\usepackage[utf8]{inputenc}
\usepackage{hyperref}
\usepackage{url}

\graphicspath{{./figures/}}

\title{CCS Concepts for LaTeX-to-DOCX Conversion Benchmark}

\author{First Author}
\affiliation{\institution{Example University}
  \city{City} \country{Country}}
\author{Second Author}
\affiliation{\institution{Another Institute}
  \city{City} \country{Country}}

% ---- ACM CCS Concepts ----
\beginCCS
\begin{textblock}{0.5}(-1.0,-0.5)
\ccsdesc[500]{Software and its engineering~Software libraries}
\ccsdesc[300]{Human-centered computing~Visualization techniques}
\ccsdesc[100]{Computing methodologies~Machine learning}
\end{textblock}
\endCCS

\begin{document}
\maketitle

\begin{abstract}
This paper presents a comprehensive benchmark for evaluating LaTeX-to-DOCX
conversion tools. We assess conversion quality across multiple dimensions including
mathematical notation, cross-references, bibliographic citations, and visual elements.
Our findings reveal significant variations in conversion accuracy across different tools
and document types.
\end{abstract}

\keywords{LaTeX, DOCX, document conversion, benchmark, evaluation}

\section{Introduction}

LaTeX is the de facto standard for scientific document preparation, particularly
in fields requiring complex mathematical notation. However, the widespread use
of Microsoft Word in academic and industry settings necessitates conversion of
LaTeX documents to DOCX format.

\section{Background}

The conversion of LaTeX to DOCX presents several technical challenges. Mathematical
equations must be converted to Office Math ML (OMML) format~\cite{lamport1994latex},
while bibliographic citations require mapping between BibTeX and Word's citation system.

\section{Experimental Setup}

We evaluate four popular conversion tools on a diverse corpus of 50 scientific documents.

\begin{table}[t]
\caption{API usage comparison across three visualization libraries.}
\label{tab:api}
\begin{tabular}{l|p{4.5cm}|p{3.5cm}}
  \toprule
  Library & Typical Call Pattern & Extensibility \\
  \midrule
  D3.js & \texttt{d3.select().data().enter().append()} & High (JS) \\
  Vega-Lite & \texttt{vl.markBar().encode(x,y).run()} & Medium (JSON) \\
  Observable Plot & \texttt{Plot.plot(\{marks: [...] \})} & Medium (JS) \\
  ECharts & \texttt{chart.setOption(\{\})} & High (JS) \\
  \bottomrule
\end{tabular}
\end{table}

\section{Results}

Our experiments reveal that mathematical formula conversion remains the most
challenging aspect, with an average accuracy of 78\% across tools.

\section{Conclusion}

This benchmark provides valuable insights for researchers and practitioners
selecting LaTeX-to-DOCX conversion tools.

\section{Acknowledgments}

This work was supported by the Example Research Foundation.

\bibliographystyle{ACM-Reference-Format}
\bibliography{refs}

\end{document}
